\section{Conclusiones}

\begin{enumerate}
    \item \textbf{Efectividad del enfoque ágil híbrido:} La combinación de ciclos cortos de Scrum, el control visual de límites WIP de Kanban y la disciplina de pruebas automatizadas continuas de Extreme Programming (XP) permitieron al equipo de tres estudiantes entregar tres incrementos completos, funcionales y robustos dentro del plazo establecido.

    \item \textbf{Arquitectura técnica escalable y segura:} La adopción de Laravel 13 junto con MySQL 8.0, combinada con transacciones pesimistas (\texttt{lockForUpdate}) y servicios desacoplados, resolvió problemas críticos de concurrencia como la sobreventa de vacantes y la duplicidad de asistencias en aula.

    \item \textbf{Innovación en certificación y actas institucionales:} La integración de firmas criptográficas HMAC-SHA-256 en los diplomas y el blindaje de inmutabilidad de las actas académicas bajo el estatuto de la UNSAAC ($\ge 75\%$ de asistencia y nota $\ge 11.00$) erradica el fraude documental y simplifica la auditoría universitaria mediante reportes oficiales A4 y comprobación pública ciudadana.

    \item \textbf{Garantía de calidad mediante pruebas continuas:} La construcción de una suite de 115 pruebas automatizadas (470 aserciones) con Pest PHP, complementada con análisis estático PHPStan Nivel 7 y chequeo de tipos estricto en TypeScript, garantiza que el software opere sin regresiones y con máxima confiabilidad en producción.

    \item \textbf{Liderazgo y uso ético de la Inteligencia Artificial:} La asistencia de herramientas de IA aceleró consultas puntuales de sintaxis y diseño criptográfico sin comprometer la autoría del equipo, manteniendo siempre el control riguroso del código mediante validaciones cruzadas entre pares.
\end{enumerate}
